% Options for packages loaded elsewhere
\PassOptionsToPackage{unicode}{hyperref}
\PassOptionsToPackage{hyphens}{url}
\documentclass[
]{article}
\usepackage{xcolor}
\usepackage{amsmath,amssymb}
\setcounter{secnumdepth}{-\maxdimen} % remove section numbering
\usepackage{iftex}
\ifPDFTeX
  \usepackage[T1]{fontenc}
  \usepackage[utf8]{inputenc}
  \usepackage{textcomp} % provide euro and other symbols
\else % if luatex or xetex
  \usepackage{unicode-math} % this also loads fontspec
  \defaultfontfeatures{Scale=MatchLowercase}
  \defaultfontfeatures[\rmfamily]{Ligatures=TeX,Scale=1}
\fi
\usepackage{lmodern}
\ifPDFTeX\else
  % xetex/luatex font selection
\fi
% Use upquote if available, for straight quotes in verbatim environments
\IfFileExists{upquote.sty}{\usepackage{upquote}}{}
\IfFileExists{microtype.sty}{% use microtype if available
  \usepackage[]{microtype}
  \UseMicrotypeSet[protrusion]{basicmath} % disable protrusion for tt fonts
}{}
\makeatletter
\@ifundefined{KOMAClassName}{% if non-KOMA class
  \IfFileExists{parskip.sty}{%
    \usepackage{parskip}
  }{% else
    \setlength{\parindent}{0pt}
    \setlength{\parskip}{6pt plus 2pt minus 1pt}}
}{% if KOMA class
  \KOMAoptions{parskip=half}}
\makeatother
\usepackage{longtable,booktabs,array}
\usepackage{caption}
\captionsetup[table]{skip=6pt}
\newcounter{none} % for unnumbered tables
\usepackage{calc} % for calculating minipage widths
% Correct order of tables after \paragraph or \subparagraph
\usepackage{etoolbox}
\makeatletter
\patchcmd\longtable{\par}{\if@noskipsec\mbox{}\fi\par}{}{}
\makeatother
\setlength{\emergencystretch}{3em} % prevent overfull lines
\providecommand{\tightlist}{%
  \setlength{\itemsep}{0pt}\setlength{\parskip}{0pt}}
\usepackage{bookmark}
\IfFileExists{xurl.sty}{\usepackage{xurl}}{} % add URL line breaks if available
\urlstyle{same}
% fallback for those not using the hyperref driver hyperxmp:
\makeatletter
\@ifundefined{xmpquote}{\newcommand{\xmpquote}[1]{#1}}{}
\makeatother
\hypersetup{
  hidelinks,
  pdfcreator={LaTeX via pandoc}}

\author{}
\date{}

\begin{document}

\section{Stage 1 Document: Requirements Specification and Project
Scope}\label{stage-1-document-requirements-specification-and-project-scope}

\textbf{Engineering Project Documentation}\\
\textbf{Faculty of Mathematics and Information Science}\\
\textbf{Warsaw University of Technology}

\begin{center}\rule{0.5\linewidth}{0.5pt}\end{center}

\subsection{1. Document Metadata and Change
History}\label{document-metadata-and-change-history}

{\def\LTcaptype{none} % do not increment counter
\begin{longtable}[]{@{}
  >{\centering\arraybackslash}p{(\linewidth - 6\tabcolsep) * \real{0.1444}}
  >{\centering\arraybackslash}p{(\linewidth - 6\tabcolsep) * \real{0.1111}}
  >{\centering\arraybackslash}p{(\linewidth - 6\tabcolsep) * \real{0.1556}}
  >{\raggedright\arraybackslash}p{(\linewidth - 6\tabcolsep) * \real{0.5889}}@{}}
\toprule\noalign{}
\begin{minipage}[b]{\linewidth}\centering
Version
\end{minipage} & \begin{minipage}[b]{\linewidth}\centering
Date
\end{minipage} & \begin{minipage}[b]{\linewidth}\centering
Author
\end{minipage} & \begin{minipage}[b]{\linewidth}\raggedright
Description of Change
\end{minipage} \\
\midrule\noalign{}
\endhead
\bottomrule\noalign{}
\endlastfoot
\textbf{1.0} & 15.10.2024 & Alice Kowalska & Baseline version for
supervisor sign-off \\
\textbf{1.1} & 18.10.2024 & Bob Nowak & Added quantifiable
non-functional requirements (NFRs) \\
\end{longtable}
}

\begin{itemize}
\tightlist
\item
  \textbf{Project Title:} \emph{A System for Reconstructing Fragmented
  Documents Based on the Shape and Content of Their Parts}
\item
  \textbf{Authors:} Martyna Sadowska, Hanna Szczerbińska, Zuzanna
  Zalewska
\item
  \textbf{Supervisor:} Janusz Rafałko, Ph.D.
\end{itemize}

\begin{center}\rule{0.5\linewidth}{0.5pt}\end{center}

\subsection{2. One-Sentence Engineering / Research Problem
Formulation}\label{one-sentence-engineering-research-problem-formulation}

\begin{quote}
\emph{``The objective of this project is to design, implement, and
evaluate a system capable of reconstructing fragmented (shredded or
torn) documents based on an image of their pieces, supported by a
dedicated document-splitting module.''}
\end{quote}

\begin{center}\rule{0.5\linewidth}{0.5pt}\end{center}

\subsection{3. Project Scope: MVP vs.~Future Work (Scope
Management)}\label{project-scope-mvp-vs.-future-work-scope-management}

\subsubsection{3.1. Minimum Viable Product (MVP -- Passing
Scope)}\label{minimum-viable-product-mvp-passing-scope}

\begin{enumerate}
\def\labelenumi{\arabic{enumi}.}
\tightlist
\item
  \textbf{Document Splitting Module (Module 1):}

  \begin{itemize}
  \tightlist
  \item
    Supports single and batch mode processing.
  \item
    Allows users to split document images using various algorithms,
    preview piece borders, re-generate if needed, and apply random
    rotations and non-overlapping background placement.
  \item
    Saves the generated fragments locally.
  \end{itemize}
\item
  \textbf{Document Reconstruction Module (Module 2):}

  \begin{itemize}
  \tightlist
  \item
    Extracts individual fragments from an input image (single or batch
    mode).
  \item
    Reconstructs original documents using user-selected or automated
    algorithms.
  \item
    Displays the reconstructed output, allows iterative improvement
    (adjusting time or algorithm), and saves the final result locally.
  \end{itemize}
\end{enumerate}

\subsubsection{3.2. Out of Scope (Future
Work)}\label{out-of-scope-future-work}

\begin{enumerate}
\def\labelenumi{\arabic{enumi}.}
\tightlist
\item
  Integration with external cloud storage providers or dedicated mobile
  applications.
\item
  Reconstruction of documents with missing fragments.
\end{enumerate}

\begin{center}\rule{0.5\linewidth}{0.5pt}\end{center}

\subsection{4. Functional Requirements (User Stories \& Acceptance
Criteria)}\label{functional-requirements-user-stories-acceptance-criteria}

\subsubsection{\texorpdfstring{\textbf{US-01: Loading a document image
for
splitting}}{US-01: Loading a document image for splitting}}\label{us-01-loading-a-document-image-for-splitting}

\begin{itemize}
\tightlist
\item
  \textbf{As a:} User of the document splitting module,
\item
  \textbf{I want to:} Upload a single document image or select a folder
  of document images,
\item
  \textbf{So that:} The system can process them for fragment generation.
\end{itemize}

\begin{quote}
\textbf{Acceptance Criteria (Gherkin):}

\begin{itemize}
\tightlist
\item
  \textbf{Given:} The splitting module is open,
\item
  \textbf{When:} The user uploads a valid image file or selects a folder
  containing valid image files,
\item
  \textbf{Then:} The system loads the images and displays a preview of
  the first document.
\item
  \textbf{Scenario (Invalid input):}
\item
  \textbf{When:} The user uploads a file in an unsupported format or an
  empty folder,
\item
  \textbf{Then:} The system displays an error message indicating the
  accepted formats.
\end{itemize}
\end{quote}

\begin{center}\rule{0.5\linewidth}{0.5pt}\end{center}

\subsubsection{\texorpdfstring{\textbf{US-02: Splitting a document into
fragments}}{US-02: Splitting a document into fragments}}\label{us-02-splitting-a-document-into-fragments}

\begin{itemize}
\tightlist
\item
  \textbf{As a:} User of the document splitting module,
\item
  \textbf{I want to:} Select a splitting algorithm and its parameters
  and preview the proposed piece borders overlaid on the document,
\item
  \textbf{So that:} I can control how the document is fragmented.
\end{itemize}

\begin{quote}
\textbf{Acceptance Criteria (Gherkin):}

\begin{itemize}
\tightlist
\item
  \textbf{Given:} A document image has been loaded,
\item
  \textbf{When:} The user selects a splitting algorithm and confirms its
  parameters,
\item
  \textbf{Then:} The system displays the proposed piece borders overlaid
  on the original image.
\item
  \textbf{When:} The user is not satisfied with the result and clicks
  ``Generate again'',
\item
  \textbf{Then:} The system generates a new set of borders using the
  same or modified parameters.
\end{itemize}
\end{quote}

\begin{center}\rule{0.5\linewidth}{0.5pt}\end{center}

\subsubsection{\texorpdfstring{\textbf{US-03: Generating a fragmented
image}}{US-03: Generating a fragmented image}}\label{us-03-generating-a-fragmented-image}

\begin{itemize}
\tightlist
\item
  \textbf{As a:} User of the document splitting module,
\item
  \textbf{I want to:} Confirm the generated borders and have the system
  randomly rotate the pieces and arrange them on a plain background
  without overlapping,
\item
  \textbf{So that:} I get an image of document fragments for testing
\end{itemize}

\begin{quote}
\textbf{Acceptance Criteria (Gherkin):}

\begin{itemize}
\tightlist
\item
  \textbf{Given:} The user has confirmed a set of piece borders,
\item
  \textbf{When:} The user sets parameters and confirms,
\item
  \textbf{Then:} The system produces an image with all fragments placed
  on the background without overlap.
\item
  \textbf{When:} The user can save the result to a chosen disk location.
\end{itemize}
\end{quote}

\begin{center}\rule{0.5\linewidth}{0.5pt}\end{center}

\subsubsection{\texorpdfstring{\textbf{US-04: Extracting fragments from
an input
image}}{US-04: Extracting fragments from an input image}}\label{us-04-extracting-fragments-from-an-input-image}

\begin{itemize}
\tightlist
\item
  \textbf{As a:} User of the document reconstruction module,
\item
  \textbf{I want to:} Upload an image containing document fragments
  (single or batch mode),
\item
  \textbf{So that:} The system extracts individual pieces for
  reconstruction.
\end{itemize}

\begin{quote}
\textbf{Acceptance Criteria (Gherkin):}

\begin{itemize}
\tightlist
\item
  \textbf{Given:} The reconstruction module is open,
\item
  \textbf{When:} The user uploads a valid image containing
  non-overlapping fragments,
\item
  \textbf{Then:} The system identifies and extracts all individual
  fragments from the image.
\item
  \textbf{Scenario (Invalid input):}
\item
  \textbf{When:} The image does not contain identifiable fragments,
\item
  \textbf{Then:} The system displays an appropriate error message.
\end{itemize}
\end{quote}

\begin{center}\rule{0.5\linewidth}{0.5pt}\end{center}

\subsubsection{\texorpdfstring{\textbf{US-05: Reconstructing a document
from
fragments}}{US-05: Reconstructing a document from fragments}}\label{us-05-reconstructing-a-document-from-fragments}

\begin{itemize}
\tightlist
\item
  \textbf{As a:} User of the document reconstruction module,
\item
  \textbf{I want to:} Select a reconstruction algorithm (or let the
  system choose automatically) and view the reconstructed document,
\item
  \textbf{So that:} I can achieve the original document content in a
  legible form.
\end{itemize}

\begin{quote}
\textbf{Acceptance Criteria (Gherkin):}

\begin{itemize}
\tightlist
\item
  \textbf{Given:} Fragments have been extracted from the input image,
\item
  \textbf{When:} The user selects an algorithm or chooses automatic
  selection, and starts reconstruction,
\item
  \textbf{Then:} The system displays the reconstructed document image.
\end{itemize}
\end{quote}

\begin{center}\rule{0.5\linewidth}{0.5pt}\end{center}

\subsubsection{\texorpdfstring{\textbf{US-06: Iterative enhancement of
reconstruction}}{US-06: Iterative enhancement of reconstruction}}\label{us-06-iterative-enhancement-of-reconstruction}

\begin{itemize}
\tightlist
\item
  \textbf{As a:} User of the document reconstruction module,
\item
  \textbf{I want to:} Review the reconstruction result and choose to
  accept it, continue the same algorithm for more time or switch to a
  different algorithm,
\item
  \textbf{So that:} I can iteratively improve the reconstruction
  quality.
\end{itemize}

\begin{quote}
\textbf{Acceptance Criteria (Gherkin):}

\begin{itemize}
\tightlist
\item
  \textbf{Given:} A reconstruction result is displayed,
\item
  \textbf{When:} The user selects ``Continue'' with the current
  algorithm,
\item
  \textbf{Then:} The system resumes processing.
\item
  \textbf{When:} The user selects ``Change algorithm'' and picks a new
  one,
\item
  \textbf{Then:} The system restarts reconstruction with the new
  algorithm.
\item
  \textbf{When:} The user selects ``Accept'',
\item
  \textbf{Then:} The system allows saving the result locally.
\end{itemize}
\end{quote}

\begin{center}\rule{0.5\linewidth}{0.5pt}\end{center}

\subsection{5. Quantifiable Non-Functional Requirements
(NFRs)}\label{quantifiable-non-functional-requirements-nfrs}

{\def\LTcaptype{none} % do not increment counter
\begin{longtable}[]{@{}
  >{\centering\arraybackslash}p{(\linewidth - 6\tabcolsep) * \real{0.2000}}
  >{\raggedright\arraybackslash}p{(\linewidth - 6\tabcolsep) * \real{0.2667}}
  >{\raggedright\arraybackslash}p{(\linewidth - 6\tabcolsep) * \real{0.2667}}
  >{\raggedright\arraybackslash}p{(\linewidth - 6\tabcolsep) * \real{0.2667}}@{}}
\toprule\noalign{}
\begin{minipage}[b]{\linewidth}\centering
ID
\end{minipage} & \begin{minipage}[b]{\linewidth}\raggedright
Category
\end{minipage} & \begin{minipage}[b]{\linewidth}\raggedright
Target Quantifiable Metric
\end{minipage} & \begin{minipage}[b]{\linewidth}\raggedright
Verification Method
\end{minipage} \\
\midrule\noalign{}
\endhead
\bottomrule\noalign{}
\endlastfoot
\textbf{NFR-01} & \textbf{Performance (Splitting)} & Splitting a single
A4-sized document image into at least 50 fragments completes under 60
seconds. & Automated benchmark script with sample dataset \\
\textbf{NFR-02} & \textbf{Performance (Reconstruction)} & Reconstruction
of a document with at least 50 fragments completes under 60 minutes. &
Automated benchmark script with sample dataset \\
\textbf{NFR-03} & \textbf{Scalability} & The system handles documents
split into up to 100 fragments without crashing. & Stress tess \\
\textbf{NFR-04} & \textbf{Usability} & A user with no prior training can
complete a full split-and-reconstruct workflow. & Usability test \\
\textbf{NFR-05} & \textbf{Error Handling} & 100\% of invalid inputs
produce a user-friendly error message. & test covering all invalid input
scenarios \\
\textbf{NFR-06} & \textbf{Availability} & The system operates fully
offline & Disconnected network test \\
\textbf{NFR-07} & \textbf{Security} & All image processing is performed
locally and no data is transmitted to external servers. & Manual
verification of network traffic \\
\end{longtable}
}

\begin{center}\rule{0.5\linewidth}{0.5pt}\end{center}

\subsection{6. Glossary}\label{glossary}

{\def\LTcaptype{none} % do not increment counter
\begin{longtable}[]{@{}
  >{\raggedright\arraybackslash}p{(\linewidth - 2\tabcolsep) * \real{0.1783}}
  >{\raggedright\arraybackslash}p{(\linewidth - 2\tabcolsep) * \real{0.8217}}@{}}
\toprule\noalign{}
\begin{minipage}[b]{\linewidth}\raggedright
Term
\end{minipage} & \begin{minipage}[b]{\linewidth}\raggedright
Definition
\end{minipage} \\
\midrule\noalign{}
\endhead
\bottomrule\noalign{}
\endlastfoot
\textbf{APD (Archiwum Prac Dyplomowych)} & University Thesis Archive;
official web portal of Warsaw University of Technology used for thesis
submission, anti-plagiarism verification, and archiving. \\
\textbf{Hypertable} & A TimescaleDB abstraction that presents
time-series data as a single continuous PostgreSQL table while
transparently partitioning it into time-based chunks under the hood. \\
\textbf{Isolation Forest} & An unsupervised machine learning algorithm
for anomaly detection that isolates observations by randomly selecting a
feature and split value, where anomalies require fewer splits to
isolate. \\
\textbf{MVP (Minimum Viable Product)} & The minimal functional scope of
the project fulfilling all foundational requirements for academic
project evaluation and defense. \\
\textbf{p95 Latency} & The 95th percentile of response time
distribution, indicating that 95\% of all incoming requests are served
in less than the specified duration. \\
\textbf{Rolling Z-Score} & A localized statistical metric representing
the number of standard deviations a given observation deviates from the
moving mean over a sliding temporal window. \\
\textbf{Telemetry} & Automated remote collection and transmission of
operational measurement data from distributed sensors or agents to a
central receiving station. \\
\end{longtable}
}

\begin{center}\rule{0.5\linewidth}{0.5pt}\end{center}

\subsection{7. Preliminary Risk Register and Role
Allocation}\label{preliminary-risk-register-and-role-allocation}

\emph{(Detailed Risk Register and RACI Matrix are defined in dedicated
documents:
\href{../management/risk-register-and-contingency.md}{\texttt{management/risk-register-and-contingency.md}}
and
\href{../management/team-charter-and-raci.md}{\texttt{management/team-charter-and-raci.md}}.)}

\begin{center}\rule{0.5\linewidth}{0.5pt}\end{center}

\subsection{8. Formal Supervisor
Acceptance}\label{formal-supervisor-acceptance}

\begin{itemize}
\tightlist
\item[$\boxtimes$]
  \textbf{Scope and MVP Sign-off:} The supervisor has reviewed the
  requirements specification and approves the MVP scope as the
  foundation for the engineering / master's degree project.
\end{itemize}

\textbf{Supervisor Signature:}
\ldots\ldots\ldots\ldots\ldots\ldots\ldots\ldots\ldots\ldots\ldots\ldots\ldots\ldots\ldots\ldots\ldots\ldots\ldots\ldots{}
\textbf{Date:} October 18, 2026

\end{document}
